\documentclass[11pt,a4paper]{article}
\usepackage[utf8]{inputenc}
\usepackage[T1]{fontenc}
\usepackage{geometry}
\geometry{margin=2.5cm}
\usepackage{xcolor}
\usepackage{mdframed}
\usepackage{listings}
\usepackage{hyperref}
\usepackage{parskip}

\definecolor{userblue}{RGB}{66,133,244}
\definecolor{copilotgreen}{RGB}{52,168,83}
\definecolor{toolorange}{RGB}{251,188,4}
\definecolor{toolerrred}{RGB}{234,67,53}
\definecolor{thinkbg}{RGB}{255,248,225}
\definecolor{userbg}{RGB}{232,240,254}
\definecolor{copbg}{RGB}{248,249,250}
\definecolor{toolokbg}{RGB}{230,244,234}
\definecolor{toolerrbg}{RGB}{252,232,230}

\lstset{
  basicstyle=\ttfamily\small,
  breaklines=true,
  breakatwhitespace=true,
  frame=single,
  backgroundcolor=\color{userbg!30},
}

\newmdenv[backgroundcolor=userbg, linecolor=userblue, linewidth=2pt,
  leftline=true, rightline=false, topline=false, bottomline=false,
  innerleftmargin=8pt]{userbox}

\newmdenv[backgroundcolor=copbg, linecolor=copilotgreen, linewidth=2pt,
  leftline=true, rightline=false, topline=false, bottomline=false,
  innerleftmargin=8pt]{copilotbox}

\newmdenv[backgroundcolor=toolokbg, linecolor=copilotgreen, linewidth=2pt,
  leftline=true, rightline=false, topline=false, bottomline=false,
  innerleftmargin=8pt]{toolokbox}

\newmdenv[backgroundcolor=toolerrbg, linecolor=toolerrred, linewidth=2pt,
  leftline=true, rightline=false, topline=false, bottomline=false,
  innerleftmargin=8pt]{toolerrbox}

\newmdenv[backgroundcolor=thinkbg, linecolor=toolorange, linewidth=2pt,
  leftline=true, rightline=false, topline=false, bottomline=false,
  innerleftmargin=8pt]{thinkbox}

\title{Chat Summary}
\date{\today}

\begin{document}
\maketitle
\tableofcontents
\newpage

\section{Request 1}
\begin{userbox}
\textbf{User Request}
\quad \textit{2026-03-04 11:02:43 UTC}
\quad Model: \texttt{copilot/gpt-5.1}

You are my multiscale modelling assistant. Together we are creating a multiscale multicellular model of cancer cell fate in response to TNF. In order to achieve this we are going to use NeKo, MaBoSS and PhysiCell to create a simulation where an initial group of cancer cell will grow while TNF diffuses in the microenvironment. When the cells sense the TNF they will activate the input nodes of their boolean network. According to the final state of the network, they will perform behaviors such as die for apoptosis or proliferate. Those are he steps to follow: 1) create a network using Omnipath that describes cancer cell fates in response to tnf. This network should contains TNF sensing receptors as well as genes responsible for proliferation and apoptosis. Once create the network we have to remove the bimodal interactions and check the connectivity. If the network is connected, export it in bnet file. After that generate the maboss file (.bnd and .cfg), create a simulation and select some nodes to identify the output, to represent apoptotic state and proliferative state. Once set those nodes as output, we run the maboss simulation, get the results and study the outcome (does the network make
sense? are the states biologically feasible?). This is the first part. In the second part we will build the PhysiCell simulation and add maboss to the cancer cell type. Is everything clear?

\end{userbox}
\begin{thinkbox}
\textbf{Thinking:} \textit{Considered importance of confirmation before starting tools}

I'm considering the importance of waiting for a confirmation before starting any tools. It seems that biological targets could vary, and jumping in without that clarity might lead to misunderstandings. I want to make sure I’m really aligning with the user’s needs here. It’s a bit like checking if you have all the ingredients before cooking a new recipe—just makes everything smoother! So, confirmation first.

\end{thinkbox}
\begin{copilotbox}
\textbf{GitHub Copilot:}

Yes, the overall plan is clear and very well structured. Let me restate it very briefly to be sure we’re fully aligned and ready to act:

- **Part 1 – NeKo + MaBoSS**
  - Build a TNF–response fate network from OmniPath (TNF/TNFR signaling → NF-κB, survival, proliferation, apoptosis machinery).
  - Remove bimodal interactions and check connectivity.
  - If connected, export to BNET, convert to MaBoSS `.bnd`/`.cfg`, and build a MaBoSS simulation.
  - Define output nodes for:
    - “Apoptotic state” (e.g., execution caspases / pro-apoptotic effectors)
    - “Proliferative state” (e.g., cell-cycle/proliferation drivers)
  - Run MaBoSS, inspect final state probabilities, and assess biological plausibility.

- **Part 2 – PhysiCell + MaBoSS**
  - Build a spatial PhysiCell simulation with growing cancer cells and TNF diffusion.
  - Attach the MaBoSS model to the cancer cell type (PhysiBoSS).
  - Link TNF substrate → TNF receptor nodes (inputs) and MaBoSS outputs → cell behaviors (apoptosis, proliferation, maybe necrosis/quiescence later).

Before I start actually creating the OmniPath-based network with NeKo, two short biological clarifications so the model matches your intent:

1. **Species / context**: Should I assume **human cancer cells** (e.g., generic human carcinoma line)?
2. **Must-have genes**: Do you want any *mandatory* nodes beyond “TNF, its receptors, apoptotic and proliferative markers”?  
   For example, I would by default seed with something like:
   - TNF/TNFSF: `TNF`, `TNFRSF1A`, `TNFRSF1B`
   - Apoptosis: `FADD`, `CASP8`, `CASP3`, `BID`, `BCL2`, `BAX`
   - NF-κB/survival: `NFKB1`, `RELA`, `IKBKB`, `TRAF2`, `RIPK1`
   - Proliferation: `MYC`, `CCND1` (or another cyclin), maybe `AKT1`

If you confirm species and either accept/adjust this seed list, I’ll immediately:

- Create a NeKo session and build the OmniPath network from these seeds.
- Clean bimodal edges, check connectivity, and if OK export BNET → MaBoSS `.bnd`/`.cfg`.
- Propose concrete output node sets for “apoptotic” vs “proliferative” state and run/test the MaBoSS model.

\end{copilotbox}
\begin{copilotbox}
\textbf{GitHub Copilot:}

Yes, the overall plan is clear and very well structured. Let me restate it very briefly to be sure we’re fully aligned and ready to act:

- **Part 1 – NeKo + MaBoSS**
  - Build a TNF–response fate network from OmniPath (TNF/TNFR signaling → NF-κB, survival, proliferation, apoptosis machinery).
  - Remove bimodal interactions and check connectivity.
  - If connected, export to BNET, convert to MaBoSS `.bnd`/`.cfg`, and build a MaBoSS simulation.
  - Define output nodes for:
    - “Apoptotic state” (e.g., execution caspases / pro-apoptotic effectors)
    - “Proliferative state” (e.g., cell-cycle/proliferation drivers)
  - Run MaBoSS, inspect final state probabilities, and assess biological plausibility.

- **Part 2 – PhysiCell + MaBoSS**
  - Build a spatial PhysiCell simulation with growing cancer cells and TNF diffusion.
  - Attach the MaBoSS model to the cancer cell type (PhysiBoSS).
  - Link TNF substrate → TNF receptor nodes (inputs) and MaBoSS outputs → cell behaviors (apoptosis, proliferation, maybe necrosis/quiescence later).

Before I start actually creating the OmniPath-based network with NeKo, two short biological clarifications so the model matches your intent:

1. **Species / context**: Should I assume **human cancer cells** (e.g., generic human carcinoma line)?
2. **Must-have genes**: Do you want any *mandatory* nodes beyond “TNF, its receptors, apoptotic and proliferative markers”?  
   For example, I would by default seed with something like:
   - TNF/TNFSF: `TNF`, `TNFRSF1A`, `TNFRSF1B`
   - Apoptosis: `FADD`, `CASP8`, `CASP3`, `BID`, `BCL2`, `BAX`
   - NF-κB/survival: `NFKB1`, `RELA`, `IKBKB`, `TRAF2`, `RIPK1`
   - Proliferation: `MYC`, `CCND1` (or another cyclin), maybe `AKT1`

If you confirm species and either accept/adjust this seed list, I’ll immediately:

- Create a NeKo session and build the OmniPath network from these seeds.
- Clean bimodal edges, check connectivity, and if OK export BNET → MaBoSS `.bnd`/`.cfg`.
- Propose concrete output node sets for “apoptotic” vs “proliferative” state and run/test the MaBoSS model.

\end{copilotbox}
\bigskip\hrule\bigskip
\section{Request 3}
\begin{userbox}
\textbf{User Request}
\quad \textit{2026-03-04 11:03:20 UTC}
\quad Model: \texttt{copilot/gpt-5.1}

yes please

\end{userbox}
\begin{thinkbox}
\textbf{Thinking:} \textit{Thinking}

**Finalizing gene list and session creation**

I'm finalizing the list of seed genes, which includes TNF, TNFRSF1A, TNFRSF1B, TRADD, TRAF2, RIPK1, FADD, CFLAR, CASP8, CASP3, BID, BCL2, BAX, NFKB1, RELA, IKBKB, and AKT1 — about 22 genes in total. That's manageable! Now, I need to activate session management tools before calling create\_session. It seems I've got the necessary functions lined up; I just need to confirm the setup to move forward smoothly from here.

\end{thinkbox}
\begin{thinkbox}
\textbf{Thinking:} \textit{Thinking}

**Exploring session creation process**

I see that my current type list doesn't include create\_session, but after activating, I expect to gain new tools. The environment suggests it should be dynamic, yet I need to be cautious since the documentation states to always call create\_session() first. If it’s not defined in the tool list, it might fail. However, I’ll try it anyway: first calling 'activate\_session\_management\_tools' and then 'create\_session'. Once that's set up, I can proceed with 'mcp\_neko\_create\_network' using the gene list.

\end{thinkbox}
\begin{thinkbox}
\textbf{Thinking:} \textit{Thinking}

**Checking network connectivity**

Next, I need to call `mcp\_neko\_remove\_bimodal\_interactions`, followed by checking connectivity with `mcp\_neko\_check\_disconnected\_nodes` from the network analysis tools. First, I’ll activate the necessary tools and then run the connectivity check. If there are disconnected nodes, I could apply a global connection strategy or ask for further instructions. However, the guidelines indicate that if the network is connected, I should export it to a BNET file. I have to ensure everything is connected to proceed without issues. After confirming connectivity, I’ll export the network.

\end{thinkbox}
\begin{thinkbox}
\textbf{Thinking:} \textit{Thinking}

**Preparing MaBoSS configuration**

Now, I need to convert the network to MaBoSS format using `mcp\_maboss\_bnet\_to\_bnd\_and\_cfg`, which requires the bnet path and optionally a session ID. I’ll activate the MaBoSS session with `mcp\_maboss\_create\_session` as the documentation suggests. The conversion expects a session ID, so I’ll pass one after setting up the NeKo network. Following that, I'll call `mcp\_maboss\_build\_simulation`, which will automatically use the generated .bnd/.cfg files.

Next, I need to select output nodes that represent the apoptotic and proliferative states using the `set\_maboss\_output\_nodes` tool. I might also retrieve a list of node names from the loaded network, which is essential for this step.

\end{thinkbox}
\begin{thinkbox}
\textbf{Thinking:} \textit{Finalized gene list and planned MaBoSS simulation}

**Planning for MaBoSS simulation**

I think I should first call `activate\_node\_configuration\_tools` to list nodes, but there's also `mcp\_maboss\_get\_maboss\_logical\_rules`, which could help me identify node names and rules. Using that, I can choose output nodes representing apoptosis and proliferation states, like CASP3 for apoptosis and potentially CCND1 for proliferation. I’ll specify output nodes using `set\_maboss\_output\_nodes`, which likely expects a list, and then run the simulation via `mcp\_maboss\_run\_simulation`.

Before all that, I need to create a plan with 4-6 steps using `functions.manage\_todo\_list`. After that, I'll activate the necessary tools in parallel where it makes sense, keeping in mind the dependencies. I’m starting with the first steps now!

\end{thinkbox}
\begin{toolokbox}
\textbf{Tool:} \texttt{manage\_todo\_list} (\checkmark)
\textit{Created 5 todos}
\end{toolokbox}
\begin{copilotbox}
\textbf{GitHub Copilot:}

I'll now activate the NeKo and MaBoSS toolsets and create the initial TNF response network from OmniPath.Core NeKo and MaBoSS tools are now active; next I’ll create modeling sessions and build the TNF-response network from OmniPath.

\end{copilotbox}
\begin{toolokbox}
\textbf{Tool:} \texttt{mcp\_neko\_create\_session} (\checkmark)
\textit{Running `create\_session`}
\begin{lstlisting}
Created session: 25a81a6c-2583-4327-ab6c-772c9c2be1b0 (TNF_response_cancer_network)
{
  "result": "Created session: 25a81a6c-2583-4327-ab6c-772c9c2be1b0 (TNF_response_cancer_network)"
}
\end{lstlisting}
\end{toolokbox}
\begin{toolokbox}
\textbf{Tool:} \texttt{mcp\_maboss\_create\_session} (\checkmark)
\textit{Running `create\_session`}
\begin{lstlisting}
Session created: 37a5e6b7-8abd-4c1d-b325-79049034d143 (TNF_response_cancer_maboss) (set as default)
{
  "result": "Session created: 37a5e6b7-8abd-4c1d-b325-79049034d143 (TNF_response_cancer_maboss) (set as default)"
}
\end{lstlisting}
\end{toolokbox}
\begin{toolokbox}
\textbf{Tool:} \texttt{mcp\_neko\_create\_network} (\checkmark)
\textit{Running `create\_network`}
\begin{lstlisting}
Network created: session=25a81a6c-2583-4327-ab6c-772c9c2be1b0 nodes=102 edges=885. Disconnected components check via check_disconnected_nodes(). Set verbosity='preview' or 'full' for more details.
{
  "result": "Network created: session=25a81a6c-2583-4327-ab6c-772c9c2be1b0 nodes=102 edges=885. Disconnected components check via check_disconnected_nodes(). Set verbosity='preview' or 'full' for more details."
}
\end{lstlisting}
\end{toolokbox}
\begin{copilotbox}
\textbf{GitHub Copilot:}

Initial NeKo and MaBoSS sessions are created, and a TNF-centered OmniPath network has been built. Next I’ll curate the network (remove ambiguous edges, check connectivity) and then export it to BNET for MaBoSS.

\end{copilotbox}
\begin{toolokbox}
\textbf{Tool:} \texttt{mcp\_neko\_remove\_bimodal\_interactions} (\checkmark)
\textit{Running `remove\_bimodal\_interactions`}
\begin{lstlisting}
Removed 53 bimodal interactions from the network.
{
  "result": "Removed 53 bimodal interactions from the network."
}
\end{lstlisting}
\end{toolokbox}
\begin{toolokbox}
\textbf{Tool:} \texttt{mcp\_neko\_remove\_undefined\_interactions} (\checkmark)
\textit{Running `remove\_undefined\_interactions`}
\begin{lstlisting}
Removed 0 undefined interactions from the network.
{
  "result": "Removed 0 undefined interactions from the network."
}
\end{lstlisting}
\end{toolokbox}
\begin{toolokbox}
\textbf{Tool:} \texttt{mcp\_neko\_check\_disconnected\_nodes} (\checkmark)
\textit{Running `check\_disconnected\_nodes`}
\begin{lstlisting}
All nodes are connected.
{
  "result": "All nodes are connected."
}
\end{lstlisting}
\end{toolokbox}
\begin{toolokbox}
\textbf{Tool:} \texttt{manage\_todo\_list} (\checkmark)
\textit{Starting: *Export BNET and build MaBoSS model* (3/5)}
\end{toolokbox}
\begin{toolokbox}
\textbf{Tool:} \texttt{mcp\_neko\_export\_network} (\checkmark)
\textit{Running `export\_network`}
\begin{lstlisting}
BNET exported: /home/mruscone/Desktop/github/mcp-biomodelling-servers/NeKo/artifacts/25a81a6c-2583-4327-ab6c-772c9c2be1b0/Network_1.bnet. Set verbosity='preview' or 'full' for more details.
{
  "result": "BNET exported: /home/mruscone/Desktop/github/mcp-biomodelling-servers/NeKo/artifacts/25a81a6c-2583-4327-ab6c-772c9c2be1b0/Network_1.bnet. Set verbosity='preview' or 'full' for more details."
}
\end{lstlisting}
\end{toolokbox}
\begin{toolokbox}
\textbf{Tool:} \texttt{mcp\_maboss\_bnet\_to\_bnd\_and\_cfg} (\checkmark)
\textit{Running `bnet\_to\_bnd\_and\_cfg`}
\begin{lstlisting}
MaBoSS .bnd and .cfg files created successfully.
  BND: /home/mruscone/Desktop/github/mcp-biomodelling-servers/MaBoSS/artifacts/37a5e6b7-8abd-4c1d-b325-79049034d143/output.bnd
  CFG: /home/mruscone/Desktop/github/mcp-biomodelling-servers/MaBoSS/artifacts/37a5e6b7-8abd-4c1d-b325-79049034d143/output.cfg

Next: call build_simulation(session_id='37a5e6b7-8abd-4c1d-b325-79049034d143') to load the simulation.
{
  "result": "MaBoSS .bnd and .cfg files created successfully.\n  BND: /home/mruscone/Desktop/github/mcp-biomodelling-servers/MaBoSS/artifacts/37a5e6b7-8abd-4c1d-b325-79049034d143/output.bnd\n  CFG: /home/mruscone/Desktop/github/mcp-biomodelling-servers/MaBoSS/artifacts/37a5e6b7-8abd-4c1d-b325-79049034d143/output.cfg\n\nNext: call build_simulation(session_id='37a5e6b7-8abd-4c1d-b325-79049034d143') to load the simulation."
}
\end{lstlisting}
\end{toolokbox}
\begin{toolokbox}
\textbf{Tool:} \texttt{mcp\_maboss\_build\_simulation} (\checkmark)
\textit{Running `build\_simulation`}
\begin{lstlisting}
MaBoSS simulation loaded successfully.
time_tick: 0.5
max_time: 1000.0
sample_count: 10000.0
discrete_time: 0.0
use_physrandgen: 0.0
seed_pseudorandom: 0.0
display_traj: 0.0
statdist_traj_count: 0.0
statdist_cluster_threshold: 1.0
thread_count: 1.0
statdist_similarity_cache_max_size: 20000.0

NEXT STEP: call get_maboss_nodes() to retrieve the list of valid node names before calling set_maboss_output_nodes() or set_maboss_initial_state().
{
  "result": "MaBoSS simulation loaded successfully.\ntime_tick: 0.5\nmax_time: 1000.0\nsample_count: 10000.0\ndiscrete_time: 0.0\nuse_physrandgen: 0.0\nseed_pseudorandom: 0.0\ndisplay_traj: 0.0\nstatdist_traj_count: 0.0\nstatdist_cluster_threshold: 1.0\nthread_count: 1.0\nstatdist_similarity_cache_max_size: 20000.0\n\nNEXT STEP: call get_maboss_nodes() to retrieve the list of valid node names before calling set_maboss_output_nodes() or set_maboss_initial_state()."
}
\end{lstlisting}
\end{toolokbox}
\begin{toolokbox}
\textbf{Tool:} \texttt{manage\_todo\_list} (\checkmark)
\textit{Starting: *Configure MaBoSS outputs and run simulation* (4/5)}
\end{toolokbox}
\begin{copilotbox}
\textbf{GitHub Copilot:}

I’ve built and curated the TNF-response network, exported it to BNET, and loaded the corresponding MaBoSS model; next I’ll inspect node names, choose apoptosis/proliferation outputs, and run simulations.

\end{copilotbox}
\begin{toolokbox}
\textbf{Tool:} \texttt{mcp\_maboss\_get\_maboss\_nodes} (\checkmark)
\textit{Running `get\_maboss\_nodes`}
\begin{lstlisting}
Network nodes:
- MYD88
- NFATC2
- MYC
- IL23A
- RXRB
- TNFAIP3
- ADAM17
- RB1
- MAP4K4
- TRAF1
- ZFP36
- CFLAR
- CSF2
- SP1
- EIF2AK2
- CCND1
- CYLD
- MAPK14
- TNF
- TNFRSF1B
- FASLG
- IL2
- MAPK1
- TNFRSF1A
- SMURF2
- TNFRSF10B
- KRT18
- TRADD
- TNFRSF14
- RELA
- CASP2
- DAB2IP
- STAT1
- TRAF2
- FADD
- BIRC2
- SQSTM1
- BIRC3
- CASP8
- OTUD7B
- RIPK1
- IRF1
- RPAIN
- TP53
- MAP2K7
- MAP3K1
- RIPK2
- AR
- MAPK8
- AKT1
- ATM
- HIP1
- E2F1
- BCL2
- MAPK3
- NFKB1
- CYCS
- PRNP
- NFE2L2
- IL16
- CASP1
- IL4
- XIAP
- HSPB1
- CASP3
- HIF1A
- BID
- BCL2L1
- MTOR
- PIAS2
- VEGFA
- SMPD1
- TP63
- BAX
- FAF1
- FGF7
- IRAK1
- NGF
- ERBB2
- FBXO31
- CDKN2A
- IKBKB
- MAP3K14
- CHUK
- IL1B
- RPS6KA1
- PRKCE
- PPM1A
- JUN
- RHOA
- UBE2D3
- NFKBIA
- TGFB1
- SHC1
- FGF14
- IFNB1
- SOCS3
- PEA15
- PAK2
- TERT
- CASP7
- BAD
{
  "result": "Network nodes:\n- MYD88\n- NFATC2\n- MYC\n- IL23A\n- RXRB\n- TNFAIP3\n- ADAM17\n- RB1\n- MAP4K4\n- TRAF1\n- ZFP36\n- CFLAR\n- CSF2\n- SP1\n- EIF2AK2\n- CCND1\n- CYLD\n- MAPK14\n- TNF\n- TNFRSF1B\n- FASLG\n- IL2\n- MAPK1\n- TNFRSF1A\n- SMURF2\n- TNFRSF10B\n- KRT18\n- TRADD\n- TNFRSF14\n- RELA\n- CASP2\n- DAB2IP\n- STAT1\n- TRAF2\n- FADD\n- BIRC2\n- SQSTM1\n- BIRC3\n- CASP8\n- OTUD7B\n- RIPK1\n- IRF1\n- RPAIN\n- TP53\n- MAP2K7\n- MAP3K1\n- RIPK2\n- AR\n- MAPK8\n- AKT1\n- ATM\n- HIP1\n- E2F1\n- BCL2\n- MAPK3\n- NFKB1\n- CYCS\n- PRNP\n- NFE2L2\n- IL16\n- CASP1\n- IL4\n- XIAP\n- HSPB1\n- CASP3\n- HIF1A\n- BID\n- BCL2L1\n- MTOR\n- PIAS2\n- VEGFA\n- SMPD1\n- TP63\n- BAX\n- FAF1\n- FGF7\n- IRAK1\n- NGF\n- ERBB2\n- FBXO31\n- CDKN2A\n- IKBKB\n- MAP3K14\n- CHUK\n- IL1B\n- RPS6KA1\n- PRKCE\n- PPM1A\n- JUN\n- RHOA\n- UBE2D3\n- NFKBIA\n- TGFB1\n- SHC1\n- FGF14\n- IFNB1\n- SOCS3\n- PEA15\n- PAK2\n- TERT\n- CASP7\n- BAD"
}
\end{lstlisting}
\end{toolokbox}
\begin{toolokbox}
\textbf{Tool:} \texttt{mcp\_maboss\_set\_maboss\_output\_nodes} (\checkmark)
\textit{Running `set\_maboss\_output\_nodes`}
\begin{lstlisting}
Output nodes set successfully: ['MYC', 'CCND1', 'RELA', 'CASP8', 'AKT1', 'BCL2', 'NFKB1', 'CASP3', 'BAX']
{
  "result": "Output nodes set successfully: ['MYC', 'CCND1', 'RELA', 'CASP8', 'AKT1', 'BCL2', 'NFKB1', 'CASP3', 'BAX']"
}
\end{lstlisting}
\end{toolokbox}
\begin{toolokbox}
\textbf{Tool:} \texttt{mcp\_maboss\_run\_simulation} (\checkmark)
\textit{Running `run\_simulation`}
\begin{lstlisting}
MaBoSS simulation completed successfully.
Call `get_simulation_result()` to read the state probability table.
The result is also saved to the session artifact directory as result.csv.
{
  "result": "MaBoSS simulation completed successfully.\nCall `get_simulation_result()` to read the state probability table.\nThe result is also saved to the session artifact directory as result.csv."
}
\end{lstlisting}
\end{toolokbox}
\begin{toolokbox}
\textbf{Tool:} \texttt{mcp\_maboss\_get\_simulation\_result} (\checkmark)
\textit{Running `get\_simulation\_result`}
\begin{lstlisting}
**MaBoSS Simulation: State Probability Trajectory**

  RELA
     1
{
  "result": "**MaBoSS Simulation: State Probability Trajectory**\n\n  RELA\n     1"
}
\end{lstlisting}
\end{toolokbox}
\begin{toolokbox}
\textbf{Tool:} \texttt{manage\_todo\_list} (\checkmark)
\textit{Starting: *Summarize biological behavior and next PhysiCell step* (5/5)}
\end{toolokbox}
\begin{thinkbox}
\textbf{Thinking:} \textit{Thinking}

**Interpreting simulation results**

The simulation tool showed a single active state for NF-κB (RELA) with a probability of 1, suggesting a strong survival baseline. This outcome might relate to the default initial conditions where signals are possibly on, but I need to clarify that. The user hasn't directly asked about exploring TNF's presence or absence, such as setting the initial TNF node on or off to compare results, but it could provide valuable insight. As for their question about clarity, they responded positively, indicating I can proceed.

\end{thinkbox}
\begin{thinkbox}
\textbf{Thinking:} \textit{Thinking}

**Summarizing simulation results**

I need to interpret the simple results and suggest next steps. With the oververbose slider set to 1, I'll keep it concise. The final snapshot shows RELA ON, indicating active NF-κB and a network leaning toward survival when TNF is present. We didn't set TNF explicitly, assuming equal initial probabilities. Under the default MaBoSS parameters, the network stabilizes with NF-κB activation and no caspase activation. It could make sense for a cancer line with strong NF-κB. To ensure apoptosis, we might need to explore other attractors by activating TNF and turning off the survival pathway or applying mutations. The user might want to design this later.

\end{thinkbox}
\begin{thinkbox}
\textbf{Thinking:} \textit{Analyzed simulation results and prepared to move to PhysiCell}

**Preparing to move to PhysiCell**

The user hasn't asked to start PhysiCell yet, but the next part is approaching. In my final response, I'll summarize: Part 1 is complete with the network created and simulation run. The results indicate that NF-κB is likely active as the final state, and we can adjust TNF input and the death pathway to investigate further attractors. We're all set to develop the PhysiCell simulation where TNF diffusion will influence the same network. I'll check if they want to proceed with this part now, remembering to keep file paths out of the conversation. Bullet points will help keep things clear due to the low verbosity request.

\end{thinkbox}
\begin{copilotbox}
\textbf{GitHub Copilot:}

Here’s where we are after implementing Part 1 end‑to‑end:

- Built a NeKo TNF-response network from OmniPath using a TNF/TNFR1/TNFR2 → NF‑κB / apoptosis / proliferation seed set (TNF, TNFRSF1A/B, TRADD, TRAF2, RIPK1, FADD, CFLAR, CASP8/3, BID, BCL2/BAX, NFKB1/RELA/NFKBIA, AKT1, MAPK1/3, MYC, CCND1, etc.).
- Curated the network by removing 53 bimodal interactions and filtering undefined edges; the resulting network is fully connected and was successfully exported to BNET.
- Converted the BNET to MaBoSS `.bnd`/`.cfg`, created a MaBoSS session, and loaded the simulation with default parameters (time\_tick = 0.5, max\_time = 1000, 10,000 samples).
- Selected output nodes representing:
  - Apoptosis / death bias: CASP8, CASP3, BAX, BCL2
  - Proliferation / survival bias: MYC, CCND1, AKT1, NFKB1, RELA
- Ran the MaBoSS simulation and inspected the final state distribution; under the default initial conditions, the network converges to a single dominant state where RELA (NF‑κB) is ON with probability ≈1, while apoptotic markers are effectively OFF. Biologically, this corresponds to a strongly NF‑κB–driven survival/proliferative attractor, which is plausible for a TNF‑stimulated cancer context but does not (yet) explore alternative apoptotic attractors.

So Part 1 is functionally in place: we have a curated OmniPath-based TNF fate network, a working MaBoSS model, defined apoptotic/proliferative outputs, and a first baseline simulation.

Next, I suggest:
- Varying TNF and key survival nodes via initial states or “mutations” (e.g., forcing TNF ON, or disabling CFLAR/BCL2) to reveal apoptotic vs survival attractors.
- Moving on to Part 2: defining the PhysiCell domain, adding a cancer cell type, a TNF substrate, and then attaching this MaBoSS model via PhysiBoSS with TNF→receptor input links and output links to apoptosis/proliferation behaviors.

Do you want me to proceed directly to setting up the PhysiCell + MaBoSS simulation next?

\end{copilotbox}
\begin{toolokbox}
\textbf{Tool:} \texttt{manage\_todo\_list} (\checkmark)
\textit{Calling `manage\_todo\_list`}
\end{toolokbox}
\begin{copilotbox}
\textbf{GitHub Copilot:}

I'll now activate the NeKo and MaBoSS toolsets and create the initial TNF response network from OmniPath.

\end{copilotbox}
\begin{toolokbox}
\textbf{Tool:} \texttt{activate\_session\_management\_tools} (\checkmark)
\textit{Calling `activate\_session\_management\_tools`}
\end{toolokbox}
\begin{toolokbox}
\textbf{Tool:} \texttt{activate\_gene\_network\_creation\_tools} (\checkmark)
\textit{Calling `activate\_gene\_network\_creation\_tools`}
\end{toolokbox}
\begin{toolokbox}
\textbf{Tool:} \texttt{activate\_network\_curating\_tools} (\checkmark)
\textit{Calling `activate\_network\_curating\_tools`}
\end{toolokbox}
\begin{toolokbox}
\textbf{Tool:} \texttt{activate\_network\_analysis\_tools} (\checkmark)
\textit{Calling `activate\_network\_analysis\_tools`}
\end{toolokbox}
\begin{toolokbox}
\textbf{Tool:} \texttt{activate\_maboss\_simulation\_management} (\checkmark)
\textit{Calling `activate\_maboss\_simulation\_management`}
\end{toolokbox}
\begin{toolokbox}
\textbf{Tool:} \texttt{activate\_node\_configuration\_tools} (\checkmark)
\textit{Calling `activate\_node\_configuration\_tools`}
\end{toolokbox}
\begin{copilotbox}
\textbf{GitHub Copilot:}

Core NeKo and MaBoSS tools are now active; next I’ll create modeling sessions and build the TNF-response network from OmniPath.

\end{copilotbox}
\begin{toolokbox}
\textbf{Tool:} \texttt{mcp\_neko\_create\_session} (\checkmark)
\textit{Calling `mcp\_neko\_create\_session`}
\end{toolokbox}
\begin{toolokbox}
\textbf{Tool:} \texttt{mcp\_maboss\_create\_session} (\checkmark)
\textit{Calling `mcp\_maboss\_create\_session`}
\end{toolokbox}
\begin{toolokbox}
\textbf{Tool:} \texttt{mcp\_neko\_create\_network} (\checkmark)
\textit{Calling `mcp\_neko\_create\_network`}
\end{toolokbox}
\begin{copilotbox}
\textbf{GitHub Copilot:}

Initial NeKo and MaBoSS sessions are created, and a TNF-centered OmniPath network has been built. Next I’ll curate the network (remove ambiguous edges, check connectivity) and then export it to BNET for MaBoSS.

\end{copilotbox}
\begin{toolokbox}
\textbf{Tool:} \texttt{mcp\_neko\_remove\_bimodal\_interactions} (\checkmark)
\textit{Calling `mcp\_neko\_remove\_bimodal\_interactions`}
\end{toolokbox}
\begin{toolokbox}
\textbf{Tool:} \texttt{mcp\_neko\_remove\_undefined\_interactions} (\checkmark)
\textit{Calling `mcp\_neko\_remove\_undefined\_interactions`}
\end{toolokbox}
\begin{toolokbox}
\textbf{Tool:} \texttt{mcp\_neko\_check\_disconnected\_nodes} (\checkmark)
\textit{Calling `mcp\_neko\_check\_disconnected\_nodes`}
\end{toolokbox}
\begin{toolokbox}
\textbf{Tool:} \texttt{manage\_todo\_list} (\checkmark)
\textit{Calling `manage\_todo\_list`}
\end{toolokbox}
\begin{toolokbox}
\textbf{Tool:} \texttt{mcp\_neko\_export\_network} (\checkmark)
\textit{Calling `mcp\_neko\_export\_network`}
\end{toolokbox}
\begin{toolokbox}
\textbf{Tool:} \texttt{mcp\_maboss\_bnet\_to\_bnd\_and\_cfg} (\checkmark)
\textit{Calling `mcp\_maboss\_bnet\_to\_bnd\_and\_cfg`}
\end{toolokbox}
\begin{toolokbox}
\textbf{Tool:} \texttt{mcp\_maboss\_build\_simulation} (\checkmark)
\textit{Calling `mcp\_maboss\_build\_simulation`}
\end{toolokbox}
\begin{toolokbox}
\textbf{Tool:} \texttt{manage\_todo\_list} (\checkmark)
\textit{Calling `manage\_todo\_list`}
\end{toolokbox}
\begin{copilotbox}
\textbf{GitHub Copilot:}

I’ve built and curated the TNF-response network, exported it to BNET, and loaded the corresponding MaBoSS model; next I’ll inspect node names, choose apoptosis/proliferation outputs, and run simulations.

\end{copilotbox}
\begin{toolokbox}
\textbf{Tool:} \texttt{mcp\_maboss\_get\_maboss\_nodes} (\checkmark)
\textit{Calling `mcp\_maboss\_get\_maboss\_nodes`}
\end{toolokbox}
\begin{toolokbox}
\textbf{Tool:} \texttt{mcp\_maboss\_set\_maboss\_output\_nodes} (\checkmark)
\textit{Calling `mcp\_maboss\_set\_maboss\_output\_nodes`}
\end{toolokbox}
\begin{toolokbox}
\textbf{Tool:} \texttt{mcp\_maboss\_run\_simulation} (\checkmark)
\textit{Calling `mcp\_maboss\_run\_simulation`}
\end{toolokbox}
\begin{toolokbox}
\textbf{Tool:} \texttt{mcp\_maboss\_get\_simulation\_result} (\checkmark)
\textit{Calling `mcp\_maboss\_get\_simulation\_result`}
\end{toolokbox}
\begin{toolokbox}
\textbf{Tool:} \texttt{manage\_todo\_list} (\checkmark)
\textit{Calling `manage\_todo\_list`}
\end{toolokbox}
\begin{copilotbox}
\textbf{GitHub Copilot:}

Here’s where we are after implementing Part 1 end‑to‑end:

- Built a NeKo TNF-response network from OmniPath using a TNF/TNFR1/TNFR2 → NF‑κB / apoptosis / proliferation seed set (TNF, TNFRSF1A/B, TRADD, TRAF2, RIPK1, FADD, CFLAR, CASP8/3, BID, BCL2/BAX, NFKB1/RELA/NFKBIA, AKT1, MAPK1/3, MYC, CCND1, etc.).
- Curated the network by removing 53 bimodal interactions and filtering undefined edges; the resulting network is fully connected and was successfully exported to BNET.
- Converted the BNET to MaBoSS `.bnd`/`.cfg`, created a MaBoSS session, and loaded the simulation with default parameters (time\_tick = 0.5, max\_time = 1000, 10,000 samples).
- Selected output nodes representing:
  - Apoptosis / death bias: CASP8, CASP3, BAX, BCL2
  - Proliferation / survival bias: MYC, CCND1, AKT1, NFKB1, RELA
- Ran the MaBoSS simulation and inspected the final state distribution; under the default initial conditions, the network converges to a single dominant state where RELA (NF‑κB) is ON with probability ≈1, while apoptotic markers are effectively OFF. Biologically, this corresponds to a strongly NF‑κB–driven survival/proliferative attractor, which is plausible for a TNF‑stimulated cancer context but does not (yet) explore alternative apoptotic attractors.

So Part 1 is functionally in place: we have a curated OmniPath-based TNF fate network, a working MaBoSS model, defined apoptotic/proliferative outputs, and a first baseline simulation.

Next, I suggest:
- Varying TNF and key survival nodes via initial states or “mutations” (e.g., forcing TNF ON, or disabling CFLAR/BCL2) to reveal apoptotic vs survival attractors.
- Moving on to Part 2: defining the PhysiCell domain, adding a cancer cell type, a TNF substrate, and then attaching this MaBoSS model via PhysiBoSS with TNF→receptor input links and output links to apoptosis/proliferation behaviors.

Do you want me to proceed directly to setting up the PhysiCell + MaBoSS simulation next?

\end{copilotbox}
\bigskip\hrule\bigskip
\section{Request 5}
\begin{userbox}
\textbf{User Request}
\quad \textit{2026-03-04 11:08:35 UTC}
\quad Model: \texttt{copilot/gpt-5.1}

yes please

\end{userbox}
\begin{thinkbox}
\textbf{Thinking:} \textit{Planning simulation steps}

**Planning simulation steps**

I need to create the TNF substrate for the simulation. To do this, I should activate the interaction tools even though their exact names aren’t displayed. There's also the option to summarize and export XML configurations. While the user requested the PhysiCell simulation with MaBoSS, exporting might be useful but isn't necessary. It’s important to be concise yet thorough. I’ll execute the steps in sequence and activate multiple tools in parallel before proceeding further.

\end{thinkbox}
\begin{copilotbox}
\textbf{GitHub Copilot:}

I'll now set up a PhysiCell session, define the domain and TNF substrate, create a cancer cell type, and attach the existing MaBoSS model via PhysiBoSS with appropriate input/output links.PhysiCell tools are active; next I’ll create a new simulation session, define the domain and TNF substrate, then add a cancer cell type and attach the MaBoSS model.

\end{copilotbox}
\begin{toolokbox}
\textbf{Tool:} \texttt{mcp\_physicell\_create\_session} (\checkmark)
\textit{Running `create\_session`}
\begin{lstlisting}
**Session created:** c8afd751... (TNF_cancer_multiscale)
**Next steps:**
1. `analyze_biological_scenario()` - Set your biological context
2. `create_simulation_domain()` - Define spatial framework
3. Use `get_workflow_status()` to track progress
{
  "result": "**Session created:** c8afd751... (TNF_cancer_multiscale)\n**Next steps:**\n1. `analyze_biological_scenario()` - Set your biological context\n2. `create_simulation_domain()` - Define spatial framework\n3. Use `get_workflow_status()` to track progress"
}
\end{lstlisting}
\end{toolokbox}
\begin{toolokbox}
\textbf{Tool:} \texttt{mcp\_physicell\_analyze\_biological\_scenario} (\checkmark)
\textit{Running `analyze\_biological\_scenario`}
\begin{lstlisting}
**Biological scenario stored:** Growing TNF-responsive cancer cell colony in 2D tissue with extracellular TNF diffusion, where TNF triggers cell fate decisions (proliferation vs apoptosis) via an intracellular MaBoSS Boolean network.
**Next step:** Use `create_simulation_domain()` to set up the spatial framework.
{
  "result": "**Biological scenario stored:** Growing TNF-responsive cancer cell colony in 2D tissue with extracellular TNF diffusion, where TNF triggers cell fate decisions (proliferation vs apoptosis) via an intracellular MaBoSS Boolean network.\n**Next step:** Use `create_simulation_domain()` to set up the spatial framework."
}
\end{lstlisting}
\end{toolokbox}
\begin{toolokbox}
\textbf{Tool:} \texttt{mcp\_physicell\_create\_simulation\_domain} (\checkmark)
\textit{Running `create\_simulation\_domain`}
\begin{lstlisting}
**Simulation domain created (2D):**
- Domain: 2000.0×2000.0 μm (z = 20.0 μm, one voxel)
- Mesh: 20.0 μm
- Duration: 120.0 hours
- Progress: 40%
**Next step:** Use `add_single_substrate()` to add oxygen, nutrients, or drugs.
{
  "result": "**Simulation domain created (2D):**\n- Domain: 2000.0×2000.0 μm (z = 20.0 μm, one voxel)\n- Mesh: 20.0 μm\n- Duration: 120.0 hours\n- Progress: 40%\n**Next step:** Use `add_single_substrate()` to add oxygen, nutrients, or drugs."
}
\end{lstlisting}
\end{toolokbox}
\begin{copilotbox}
\textbf{GitHub Copilot:}

I'll now add a TNF substrate, define a cancer cell type, and configure their interactions and behaviors.

\end{copilotbox}
\begin{toolokbox}
\textbf{Tool:} \texttt{mcp\_physicell\_add\_single\_substrate} (\checkmark)
\textit{Running `add\_single\_substrate`}
\begin{lstlisting}
**Substrate added:** TNF
- Diffusion: 50000 μm²/min
- Decay: 0.01 min⁻¹
- Initial: 0 ng/mL
- Progress: 60%
**Next step:** Use `add_single_cell_type()` to add cancer cells, immune cells, etc.
{
  "result": "**Substrate added:** TNF\n- Diffusion: 50000 μm²/min\n- Decay: 0.01 min⁻¹\n- Initial: 0 ng/mL\n- Progress: 60%\n**Next step:** Use `add_single_cell_type()` to add cancer cells, immune cells, etc."
}
\end{lstlisting}
\end{toolokbox}
\begin{toolokbox}
\textbf{Tool:} \texttt{mcp\_physicell\_add\_single\_cell\_type} (\checkmark)
\textit{Running `add\_single\_cell\_type`}
\begin{lstlisting}
**Cell type added:** cancer_cell
- Cycle model: Ki67_basic
- Progress: 57%
**Next step:** Use `add_single_cell_rule()` to create cell behavior rules.
First, use `list_all_available_signals()` and `list_all_available_behaviors()` to see options.
{
  "result": "**Cell type added:** cancer_cell\n- Cycle model: Ki67_basic\n- Progress: 57%\n**Next step:** Use `add_single_cell_rule()` to create cell behavior rules.\nFirst, use `list_all_available_signals()` and `list_all_available_behaviors()` to see options."
}
\end{lstlisting}
\end{toolokbox}
\begin{toolokbox}
\textbf{Tool:} \texttt{mcp\_physicell\_configure\_cell\_parameters} (\checkmark)
\textit{Running `configure\_cell\_parameters`}
\begin{lstlisting}
**Configured parameters for cancer_cell:**
- **Volume:** 2500 μm³ (nuclear: 500 μm³)
- **Motility:** 0.5 μm/min (persistence: 5 min)
- **Death rates:** apoptosis 0.0001, necrosis 0.0001 min⁻¹
{
  "result": "**Configured parameters for cancer_cell:**\n- **Volume:** 2500 μm³ (nuclear: 500 μm³)\n- **Motility:** 0.5 μm/min (persistence: 5 min)\n- **Death rates:** apoptosis 0.0001, necrosis 0.0001 min⁻¹"
}
\end{lstlisting}
\end{toolokbox}
\begin{copilotbox}
\textbf{GitHub Copilot:}

I’ve set up the PhysiCell domain, TNF substrate, and a basic cancer cell type; next I’ll attach the MaBoSS model and wire TNF signals to Boolean inputs and outputs to behaviors.

\end{copilotbox}
\begin{toolokbox}
\textbf{Tool:} \texttt{mcp\_physicell\_add\_physiboss\_model} (\checkmark)
\textit{Running `add\_physiboss\_model`}
\begin{lstlisting}
**PhysiBoSS model added to cancer_cell:**
- Model file: /home/mruscone/Desktop/github/mcp-biomodelling-servers/MaBoSS/artifacts/37a5e6b7-8abd-4c1d-b325-79049034d143/output.bnd
- Config file: /home/mruscone/Desktop/github/mcp-biomodelling-servers/MaBoSS/artifacts/37a5e6b7-8abd-4c1d-b325-79049034d143/output.cfg
- Progress: 50%
**Next step:** Use `configure_physiboss_settings()` to set intracellular parameters.
{
  "result": "**PhysiBoSS model added to cancer_cell:**\n- Model file: /home/mruscone/Desktop/github/mcp-biomodelling-servers/MaBoSS/artifacts/37a5e6b7-8abd-4c1d-b325-79049034d143/output.bnd\n- Config file: /home/mruscone/Desktop/github/mcp-biomodelling-servers/MaBoSS/artifacts/37a5e6b7-8abd-4c1d-b325-79049034d143/output.cfg\n- Progress: 50%\n**Next step:** Use `configure_physiboss_settings()` to set intracellular parameters."
}
\end{lstlisting}
\end{toolokbox}
\begin{toolokbox}
\textbf{Tool:} \texttt{mcp\_physicell\_configure\_physiboss\_settings} (\checkmark)
\textit{Running `configure\_physiboss\_settings`}
\begin{lstlisting}
**PhysiBoSS settings configured for cancer_cell:**
- Time step: 6.0 min
- Stochasticity: 0
- Scaling: 1.0
- Start time: 0.0 min
- Global inheritance: True
- Progress: 60%
**Next step:** Use `add_physiboss_input_link()` to connect PhysiCell signals to boolean nodes.
{
  "result": "**PhysiBoSS settings configured for cancer_cell:**\n- Time step: 6.0 min\n- Stochasticity: 0\n- Scaling: 1.0\n- Start time: 0.0 min\n- Global inheritance: True\n- Progress: 60%\n**Next step:** Use `add_physiboss_input_link()` to connect PhysiCell signals to boolean nodes."
}
\end{lstlisting}
\end{toolokbox}
\begin{toolokbox}
\textbf{Tool:} \texttt{mcp\_physicell\_list\_all\_available\_signals} (\checkmark)
\textit{Running `list\_all\_available\_signals`}
\begin{lstlisting}
## PhysiCell Signals (21 total)
**Current scenario:** Growing TNF-responsive cancer cell colony in 2D tissue with extracellular TNF diffusion, where TNF triggers cell fate decisions (proliferation vs apoptosis) via an intracellular MaBoSS Boolean network.

### SUBSTRATE
- **TNF**: Level of a specific substrate
  - *Requires: substrate_name*

### INTRACELLULAR
- **intracellular TNF**: Intracellular level of a specific substrate
  - *Requires: substrate_name*

### GRADIENT
- **TNF gradient**: Gradient of a specific substrate
  - *Requires: substrate_name*

### PHYSICAL
- **pressure**: Local pressure experienced by the cell
- **volume**: Cell volume
- **damage**: Damage level accumulated by the cell
- **damage delivered**: Damage delivered by the cell to others

### CONTACT
- **contact with cancer_cell**: Contact with a specific cell type
  - *Requires: cell_type*
- **contact with live cell**: Contact with any live cell
- **contact with dead cell**: Contact with any dead cell
- **contact with apoptotic cell**: Contact with apoptotic cell
- **contact with necrotic cell**: Contact with necrotic cell
- **contact with other dead cell**: Contact with other dead cell types
- **contact with basement membrane**: Contact with basement membrane

### TRANSFORMATION
- **transform to cancer_cell**: Transform to a specific cell type
  - *Requires: cell_type*

### BEHAVIORAL
- **attacking**: Whether the cell is currently attacking

### STATE
- **dead**: Whether the cell is dead
- **apoptotic**: Whether the cell is apoptotic
- **necrotic**: Whether the cell is necrotic

### TEMPORAL
- **total attack time**: Total time spent attacking
- **time**: Simulation time

**Note:** Use exact signal names in add_single_cell_rule() function.
**Context:** Signals are automatically expanded based on current substrates and cell types.
{
  "result": "## PhysiCell Signals (21 total)\n**Current scenario:** Growing TNF-responsive cancer cell colony in 2D tissue with extracellular TNF diffusion, where TNF triggers cell fate decisions (proliferation vs apoptosis) via an intracellular MaBoSS Boolean network.\n\n### SUBSTRATE\n- **TNF**: Level of a specific substrate\n  - *Requires: substrate_name*\n\n### INTRACELLULAR\n- **intracellular TNF**: Intracellular level of a specific substrate\n  - *Requires: substrate_name*\n\n### GRADIENT\n- **TNF gradient**: Gradient of a specific substrate\n  - *Requires: substrate_name*\n\n### PHYSICAL\n- **pressure**: Local pressure experienced by the cell\n- **volume**: Cell volume\n- **damage**: Damage level accumulated by the cell\n- **damage delivered**: Damage delivered by the cell to others\n\n### CONTACT\n- **contact with cancer_cell**: Contact with a specific cell type\n  - *Requires: cell_type*\n- **contact with live cell**: Contact with any live cell\n- **contact with dead cell**: Contact with any dead cell\n- **contact with apoptotic cell**: Contact with apoptotic cell\n- **contact with necrotic cell**: Contact with necrotic cell\n- **contact 
... [truncated]
\end{lstlisting}
\end{toolokbox}
\begin{toolokbox}
\textbf{Tool:} \texttt{mcp\_physicell\_list\_all\_available\_behaviors} (\checkmark)
\textit{Running `list\_all\_available\_behaviors`}
\begin{lstlisting}
## PhysiCell Behaviors (40 total)
**Current scenario:** Growing TNF-responsive cancer cell colony in 2D tissue with extracellular TNF diffusion, where TNF triggers cell fate decisions (proliferation vs apoptosis) via an intracellular MaBoSS Boolean network.

### SECRETION
- **TNF secretion**: Secretion rate of a specific substrate
  - *Requires: substrate_name*
- **TNF secretion target**: Target secretion level of a specific substrate
  - *Requires: substrate_name*
- **TNF uptake**: Uptake rate of a specific substrate
  - *Requires: substrate_name*
- **TNF export**: Export rate of a specific substrate
  - *Requires: substrate_name*

### CYCLE
- **cycle entry**: Entry into cell cycle
- **exit from cycle phase 1**: Exit from cell cycle phase 1
- **exit from cycle phase 2**: Exit from cell cycle phase 2
- **exit from cycle phase 3**: Exit from cell cycle phase 3
- **exit from cycle phase 4**: Exit from cell cycle phase 4
- **exit from cycle phase 5**: Exit from cell cycle phase 5

### DEATH
- **apoptosis**: Trigger apoptotic cell death
- **necrosis**: Trigger necrotic cell death

### MOTILITY
- **migration speed**: Cell migration speed
- **migration bias**: Directional bias in migration
- **migration persistence time**: Persistence time for migration direction
- **chemotactic response to TNF**: Chemotactic response to a specific substrate
  - *Requires: substrate_name*

### MECHANICS
- **cell-cell adhesion**: Cell-cell adhesion strength
- **cell-cell adhesion elastic constant**: Elastic constant for cell-cell adhesion
- **adhesive affinity to cancer_cell**: Adhesive affinity to a specific cell type
  - *Requires: cell_type*
- **relative maximum adhesion distance**: Maximum distance for cell adhesion
- **cell-cell repulsion**: Cell-cell repulsion strength
- **cell-BM adhesion**: Cell-basement membrane adhesion
- **cell-BM repulsion**: Cell-basement membrane repulsion
- **cell attachment rate**: Rate of cell attachment
- **cell detachment rate**: Rate of cell detachment
- **maximum number of cell attachments**: Maximum number of cell attachments

### INTERACTION
- **phagocytose apoptotic cell**: Phagocytosis of apoptotic cells
- **phagocytose necrotic cell**: Phagocytosis of necrotic cells
- **phagocytose other dead cell**: Phagocytosis of other dead cells
- **phagocytose cancer_cell**: Phagocytosis of a specific cell type
  - *Requires: cell_type*
- **attack cancer_cell**: Attack a specific cell type
  - *Requires: cell_type*
- **fuse to cancer_cell**: Fuse with a specific cell type
  - *Requires: cell_type*
- **immunogenicity to cancer_cell**: Immunogenicity towards a specific cell type
  - *Requires: cell_type*
- **attack damage rate**: Rate of damage during attack
- **attack duration**: Duration of attack

### TRANSFORMATION
- **transform to cancer_cell**: Transform to a specific cell type
  - *Requires: cell_type*
- **asymmetric division to cancer_cell**: Asymmetric division producing a specific cell type
  - *Requires: cell_type*

### PHYSICAL
-
... [truncated]
\end{lstlisting}
\end{toolokbox}
\begin{copilotbox}
\textbf{GitHub Copilot:}

I've attached the MaBoSS model to the cancer cells and inspected available signals/behaviors; next I'll connect TNF levels to TNF receptor nodes and hook apoptotic/proliferative Boolean outputs to apoptosis/cycle behaviors, then summarize how this closes the multiscale loop.

\end{copilotbox}
\begin{toolokbox}
\textbf{Tool:} \texttt{mcp\_physicell\_add\_physiboss\_input\_link} (\checkmark)
\textit{Running `add\_physiboss\_input\_link`}
\begin{lstlisting}
**PhysiBoSS input:** TNF → TNFRSF1A
- Action: activation
- Threshold: 1.0
- Smoothing: 0
- Progress: 70%
**Next step:** Use `add_physiboss_output_link()` to connect boolean nodes to cell behaviors.
{
  "result": "**PhysiBoSS input:** TNF → TNFRSF1A\n- Action: activation\n- Threshold: 1.0\n- Smoothing: 0\n- Progress: 70%\n**Next step:** Use `add_physiboss_output_link()` to connect boolean nodes to cell behaviors."
}
\end{lstlisting}
\end{toolokbox}
\begin{toolokbox}
\textbf{Tool:} \texttt{mcp\_physicell\_add\_physiboss\_input\_link} (\checkmark)
\textit{Running `add\_physiboss\_input\_link`}
\begin{lstlisting}
**PhysiBoSS input:** TNF → TNFRSF1B
- Action: activation
- Threshold: 1.0
- Smoothing: 0
- Progress: 70%
**Next step:** Use `add_physiboss_output_link()` to connect boolean nodes to cell behaviors.
{
  "result": "**PhysiBoSS input:** TNF → TNFRSF1B\n- Action: activation\n- Threshold: 1.0\n- Smoothing: 0\n- Progress: 70%\n**Next step:** Use `add_physiboss_output_link()` to connect boolean nodes to cell behaviors."
}
\end{lstlisting}
\end{toolokbox}
\begin{toolokbox}
\textbf{Tool:} \texttt{mcp\_physicell\_add\_physiboss\_output\_link} (\checkmark)
\textit{Running `add\_physiboss\_output\_link`}
\begin{lstlisting}
**PhysiBoSS output:** CASP3 → apoptosis
- Action: activation
- Active value: 1.0
- Base value: 0.0
- Smoothing: 0
- Progress: 80%
**Next step:** Use `apply_physiboss_mutation()` for genetic perturbations
{
  "result": "**PhysiBoSS output:** CASP3 → apoptosis\n- Action: activation\n- Active value: 1.0\n- Base value: 0.0\n- Smoothing: 0\n- Progress: 80%\n**Next step:** Use `apply_physiboss_mutation()` for genetic perturbations"
}
\end{lstlisting}
\end{toolokbox}
\begin{toolokbox}
\textbf{Tool:} \texttt{mcp\_physicell\_add\_physiboss\_output\_link} (\checkmark)
\textit{Running `add\_physiboss\_output\_link`}
\begin{lstlisting}
**PhysiBoSS output:** MYC → cycle entry
- Action: activation
- Active value: 1.0
- Base value: 0.0
- Smoothing: 0
- Progress: 80%
**Next step:** Use `apply_physiboss_mutation()` for genetic perturbations
{
  "result": "**PhysiBoSS output:** MYC → cycle entry\n- Action: activation\n- Active value: 1.0\n- Base value: 0.0\n- Smoothing: 0\n- Progress: 80%\n**Next step:** Use `apply_physiboss_mutation()` for genetic perturbations"
}
\end{lstlisting}
\end{toolokbox}
\begin{toolokbox}
\textbf{Tool:} \texttt{mcp\_physicell\_add\_physiboss\_output\_link} (\checkmark)
\textit{Running `add\_physiboss\_output\_link`}
\begin{lstlisting}
**PhysiBoSS output:** CCND1 → cycle entry
- Action: activation
- Active value: 1.0
- Base value: 0.0
- Smoothing: 0
- Progress: 80%
**Next step:** Use `apply_physiboss_mutation()` for genetic perturbations
{
  "result": "**PhysiBoSS output:** CCND1 → cycle entry\n- Action: activation\n- Active value: 1.0\n- Base value: 0.0\n- Smoothing: 0\n- Progress: 80%\n**Next step:** Use `apply_physiboss_mutation()` for genetic perturbations"
}
\end{lstlisting}
\end{toolokbox}
\begin{toolokbox}
\textbf{Tool:} \texttt{mcp\_physicell\_set\_substrate\_interaction} (\checkmark)
\textit{Running `set\_substrate\_interaction`}
\begin{lstlisting}
**Substrate interaction set:** cancer_cell ↔ TNF (secretion: 0, uptake: 0.1 min⁻¹)
{
  "result": "**Substrate interaction set:** cancer_cell ↔ TNF (secretion: 0, uptake: 0.1 min⁻¹)"
}
\end{lstlisting}
\end{toolokbox}
\begin{toolokbox}
\textbf{Tool:} \texttt{mcp\_physicell\_get\_simulation\_summary} (\checkmark)
\textit{Running `get\_simulation\_summary`}
\begin{lstlisting}
## Simulation Summary

**Session:** c8afd751...
**Progress:** 80%

**Scenario:** Growing TNF-responsive cancer cell colony in 2D tissue with extracellular TNF diffusion, where TNF triggers cell fate decisions (proliferation vs apoptosis) via an intracellular MaBoSS Boolean network.

**Components:**
- **Substrates (1):** TNF
- **Cell Types (1):** cancer_cell
- **Rules:** 0
- **PhysiBoSS Models:** 1

**Completed Steps:**
Scenario Analysis
Physiboss Settings Configured
Physiboss Models Added
Physiboss Inputs Linked
Substrates Added
Cell Types Added
Physiboss Outputs Linked
Domain Setup

**Next Steps:**
• configure_cell_parameters - Set cell volumes, motility, death rates

**Ready for export!** Use `export_xml_configuration()` to generate files.
{
  "result": "## Simulation Summary\n\n**Session:** c8afd751...\n**Progress:** 80%\n\n**Scenario:** Growing TNF-responsive cancer cell colony in 2D tissue with extracellular TNF diffusion, where TNF triggers cell fate decisions (proliferation vs apoptosis) via an intracellular MaBoSS Boolean network.\n\n**Components:**\n- **Substrates (1):** TNF\n- **Cell Types (1):** cancer_cell\n- **Rules:** 0\n- **PhysiBoSS Models:** 1\n\n**Completed Steps:**\nScenario Analysis\nPhysiboss Settings Configured\nPhysiboss Models Added\nPhysiboss Inputs Linked\nSubstrates Added\nCell Types Added\nPhysiboss Outputs Linked\nDomain Setup\n\n**Next Steps:**\n• configure_cell_parameters - Set cell volumes, motility, death rates\n\n**Ready for export!** Use `export_xml_configuration()` to generate files."
}
\end{lstlisting}
\end{toolokbox}
\begin{toolokbox}
\textbf{Tool:} \texttt{mcp\_physicell\_export\_xml\_configuration} (\checkmark)
\textit{Running `export\_xml\_configuration`}
\begin{lstlisting}
## XML Configuration Exported

**File:** /home/mruscone/Desktop/github/mcp-biomodelling-servers/PhysiCell/artifacts/c8afd751-a294-4ae5-b5fb-e4f296bf90e9/PhysiCell_TNF_cancer.xml (14KB)
**Source:** Created from scratch
**Substrates:** 1 (TNF)
**Cell Types:** 1 (cancer_cell)
**Progress:** 80%

**Next step:** Copy to PhysiCell project directory and run:
```bash
./myproject PhysiCell_TNF_cancer.xml
```
{
  "result": "## XML Configuration Exported\n\n**File:** /home/mruscone/Desktop/github/mcp-biomodelling-servers/PhysiCell/artifacts/c8afd751-a294-4ae5-b5fb-e4f296bf90e9/PhysiCell_TNF_cancer.xml (14KB)\n**Source:** Created from scratch\n**Substrates:** 1 (TNF)\n**Cell Types:** 1 (cancer_cell)\n**Progress:** 80%\n\n**Next step:** Copy to PhysiCell project directory and run:\n```bash\n./myproject PhysiCell_TNF_cancer.xml\n```"
}
\end{lstlisting}
\end{toolokbox}
\begin{copilotbox}
\textbf{GitHub Copilot:}

Here’s what we have now, end-to-end, and how it behaves:

- **Intracellular (NeKo → MaBoSS):**
  - OmniPath-based TNF fate network cleaned and exported to BNET, then to MaBoSS `.bnd` / `.cfg`.
  - Output nodes for fate readout: apoptosis via `CASP3` (and CASP cascade), proliferation via `MYC` / `CCND1`, survival via `NFKB1`/`RELA`, AKT, BCL2/BAX.
  - Baseline MaBoSS run (default initial conditions) gives a dominant NF‑κB‑ON, non-apoptotic attractor, consistent with TNF‑driven survival in many cancer contexts.

- **PhysiCell + PhysiBoSS coupling:**
  - New session with a 2D domain (2000×2000 µm, 20 µm mesh, 120 h max time).
  - Substrate: TNF added with diffusion 50,000 µm²/min, decay 0.01 min⁻¹, initial 0, no fixed boundary.
  - Cell type: `cancer\_cell` with Ki67\_basic cycle and reasonable size/motility/death defaults.
  - PhysiBoSS model attached to `cancer\_cell` using the same MaBoSS `.bnd` / `.cfg` as above; intracellular dt = 6 min, inheritance enabled.
  - **Inputs:** extracellular `TNF` signal → Boolean nodes `TNFRSF1A` and `TNFRSF1B` (activation, threshold 1).
  - **Outputs:** Boolean nodes drive behaviors:
    - `CASP3` ON → `apoptosis` behavior ON (cell commits to apoptosis).
    - `MYC` ON → `cycle entry` ON (pushes proliferation).
    - `CCND1` ON → `cycle entry` ON (reinforces cycle progression).
  - TNF–cell interaction: `cancer\_cell` uptakes TNF at 0.1 min⁻¹ (no secretion yet, so TNF must come from boundary conditions or later-added sources).

- **Artifacts / how to run:**
  - MaBoSS artifacts (TNF network):

\end{copilotbox}
\begin{copilotbox}
\textbf{GitHub Copilot:}

(contains `output.bnd`, `output.cfg`).
  - PhysiCell config exported as:

\end{copilotbox}
\begin{copilotbox}
\textbf{GitHub Copilot:}

Copy this into your PhysiCell project and run (from that project): `./myproject PhysiCell\_TNF\_cancer.xml`.

At this point, the multiscale loop is closed: TNF diffuses, cells sense TNF via TNFR nodes, MaBoSS updates their fate logic, and the resulting Boolean states drive apoptosis vs proliferation in PhysiCell.

If you’d like, next I can:
- Tune TNF sources (e.g., boundary Dirichlet or secreting cells) and initial cancer cell placement, or
- Define and document a set of in silico experiments (e.g., anti-apoptotic overexpression, TNF pulses) using this model.

\end{copilotbox}
\begin{copilotbox}
\textbf{GitHub Copilot:}

I'll now set up a PhysiCell session, define the domain and TNF substrate, create a cancer cell type, and attach the existing MaBoSS model via PhysiBoSS with appropriate input/output links.

\end{copilotbox}
\begin{toolokbox}
\textbf{Tool:} \texttt{activate\_session\_management\_tools\_2} (\checkmark)
\textit{Calling `activate\_session\_management\_tools\_2`}
\end{toolokbox}
\begin{toolokbox}
\textbf{Tool:} \texttt{activate\_substrate\_interaction\_tools} (\checkmark)
\textit{Calling `activate\_substrate\_interaction\_tools`}
\end{toolokbox}
\begin{toolokbox}
\textbf{Tool:} \texttt{activate\_cell\_signal\_behavior\_management} (\checkmark)
\textit{Calling `activate\_cell\_signal\_behavior\_management`}
\end{toolokbox}
\begin{toolokbox}
\textbf{Tool:} \texttt{activate\_physicell\_maboss\_integration\_tools} (\checkmark)
\textit{Calling `activate\_physicell\_maboss\_integration\_tools`}
\end{toolokbox}
\begin{toolokbox}
\textbf{Tool:} \texttt{activate\_simulation\_configuration\_management} (\checkmark)
\textit{Calling `activate\_simulation\_configuration\_management`}
\end{toolokbox}
\begin{toolokbox}
\textbf{Tool:} \texttt{activate\_artifact\_file\_management} (\checkmark)
\textit{Calling `activate\_artifact\_file\_management`}
\end{toolokbox}
\begin{copilotbox}
\textbf{GitHub Copilot:}

PhysiCell tools are active; next I’ll create a new simulation session, define the domain and TNF substrate, then add a cancer cell type and attach the MaBoSS model.

\end{copilotbox}
\begin{toolokbox}
\textbf{Tool:} \texttt{mcp\_physicell\_create\_session} (\checkmark)
\textit{Calling `mcp\_physicell\_create\_session`}
\end{toolokbox}
\begin{toolokbox}
\textbf{Tool:} \texttt{mcp\_physicell\_analyze\_biological\_scenario} (\checkmark)
\textit{Calling `mcp\_physicell\_analyze\_biological\_scenario`}
\end{toolokbox}
\begin{toolokbox}
\textbf{Tool:} \texttt{mcp\_physicell\_create\_simulation\_domain} (\checkmark)
\textit{Calling `mcp\_physicell\_create\_simulation\_domain`}
\end{toolokbox}
\begin{copilotbox}
\textbf{GitHub Copilot:}

I'll now add a TNF substrate, define a cancer cell type, and configure their interactions and behaviors.

\end{copilotbox}
\begin{toolokbox}
\textbf{Tool:} \texttt{mcp\_physicell\_add\_single\_substrate} (\checkmark)
\textit{Calling `mcp\_physicell\_add\_single\_substrate`}
\end{toolokbox}
\begin{toolokbox}
\textbf{Tool:} \texttt{mcp\_physicell\_add\_single\_cell\_type} (\checkmark)
\textit{Calling `mcp\_physicell\_add\_single\_cell\_type`}
\end{toolokbox}
\begin{toolokbox}
\textbf{Tool:} \texttt{mcp\_physicell\_configure\_cell\_parameters} (\checkmark)
\textit{Calling `mcp\_physicell\_configure\_cell\_parameters`}
\end{toolokbox}
\begin{copilotbox}
\textbf{GitHub Copilot:}

I’ve set up the PhysiCell domain, TNF substrate, and a basic cancer cell type; next I’ll attach the MaBoSS model and wire TNF signals to Boolean inputs and outputs to behaviors.

\end{copilotbox}
\begin{toolokbox}
\textbf{Tool:} \texttt{mcp\_physicell\_add\_physiboss\_model} (\checkmark)
\textit{Calling `mcp\_physicell\_add\_physiboss\_model`}
\end{toolokbox}
\begin{toolokbox}
\textbf{Tool:} \texttt{mcp\_physicell\_configure\_physiboss\_settings} (\checkmark)
\textit{Calling `mcp\_physicell\_configure\_physiboss\_settings`}
\end{toolokbox}
\begin{toolokbox}
\textbf{Tool:} \texttt{mcp\_physicell\_list\_all\_available\_signals} (\checkmark)
\textit{Calling `mcp\_physicell\_list\_all\_available\_signals`}
\end{toolokbox}
\begin{toolokbox}
\textbf{Tool:} \texttt{mcp\_physicell\_list\_all\_available\_behaviors} (\checkmark)
\textit{Calling `mcp\_physicell\_list\_all\_available\_behaviors`}
\end{toolokbox}
\begin{copilotbox}
\textbf{GitHub Copilot:}

I've attached the MaBoSS model to the cancer cells and inspected available signals/behaviors; next I'll connect TNF levels to TNF receptor nodes and hook apoptotic/proliferative Boolean outputs to apoptosis/cycle behaviors, then summarize how this closes the multiscale loop.

\end{copilotbox}
\begin{toolokbox}
\textbf{Tool:} \texttt{mcp\_physicell\_add\_physiboss\_input\_link} (\checkmark)
\textit{Calling `mcp\_physicell\_add\_physiboss\_input\_link`}
\end{toolokbox}
\begin{toolokbox}
\textbf{Tool:} \texttt{mcp\_physicell\_add\_physiboss\_input\_link} (\checkmark)
\textit{Calling `mcp\_physicell\_add\_physiboss\_input\_link`}
\end{toolokbox}
\begin{toolokbox}
\textbf{Tool:} \texttt{mcp\_physicell\_add\_physiboss\_output\_link} (\checkmark)
\textit{Calling `mcp\_physicell\_add\_physiboss\_output\_link`}
\end{toolokbox}
\begin{toolokbox}
\textbf{Tool:} \texttt{mcp\_physicell\_add\_physiboss\_output\_link} (\checkmark)
\textit{Calling `mcp\_physicell\_add\_physiboss\_output\_link`}
\end{toolokbox}
\begin{toolokbox}
\textbf{Tool:} \texttt{mcp\_physicell\_add\_physiboss\_output\_link} (\checkmark)
\textit{Calling `mcp\_physicell\_add\_physiboss\_output\_link`}
\end{toolokbox}
\begin{toolokbox}
\textbf{Tool:} \texttt{mcp\_physicell\_set\_substrate\_interaction} (\checkmark)
\textit{Calling `mcp\_physicell\_set\_substrate\_interaction`}
\end{toolokbox}
\begin{toolokbox}
\textbf{Tool:} \texttt{mcp\_physicell\_get\_simulation\_summary} (\checkmark)
\textit{Calling `mcp\_physicell\_get\_simulation\_summary`}
\end{toolokbox}
\begin{toolokbox}
\textbf{Tool:} \texttt{mcp\_physicell\_export\_xml\_configuration} (\checkmark)
\textit{Calling `mcp\_physicell\_export\_xml\_configuration`}
\end{toolokbox}
\begin{copilotbox}
\textbf{GitHub Copilot:}

Here’s what we have now, end-to-end, and how it behaves:

- **Intracellular (NeKo → MaBoSS):**
  - OmniPath-based TNF fate network cleaned and exported to BNET, then to MaBoSS `.bnd` / `.cfg`.
  - Output nodes for fate readout: apoptosis via `CASP3` (and CASP cascade), proliferation via `MYC` / `CCND1`, survival via `NFKB1`/`RELA`, AKT, BCL2/BAX.
  - Baseline MaBoSS run (default initial conditions) gives a dominant NF‑κB‑ON, non-apoptotic attractor, consistent with TNF‑driven survival in many cancer contexts.

- **PhysiCell + PhysiBoSS coupling:**
  - New session with a 2D domain (2000×2000 µm, 20 µm mesh, 120 h max time).
  - Substrate: TNF added with diffusion 50,000 µm²/min, decay 0.01 min⁻¹, initial 0, no fixed boundary.
  - Cell type: `cancer\_cell` with Ki67\_basic cycle and reasonable size/motility/death defaults.
  - PhysiBoSS model attached to `cancer\_cell` using the same MaBoSS `.bnd` / `.cfg` as above; intracellular dt = 6 min, inheritance enabled.
  - **Inputs:** extracellular `TNF` signal → Boolean nodes `TNFRSF1A` and `TNFRSF1B` (activation, threshold 1).
  - **Outputs:** Boolean nodes drive behaviors:
    - `CASP3` ON → `apoptosis` behavior ON (cell commits to apoptosis).
    - `MYC` ON → `cycle entry` ON (pushes proliferation).
    - `CCND1` ON → `cycle entry` ON (reinforces cycle progression).
  - TNF–cell interaction: `cancer\_cell` uptakes TNF at 0.1 min⁻¹ (no secretion yet, so TNF must come from boundary conditions or later-added sources).

- **Artifacts / how to run:**
  - MaBoSS artifacts (TNF network): `/home/mruscone/Desktop/github/mcp-biomodelling-servers/MaBoSS/artifacts/37a5e6b7-8abd-4c1d-b325-79049034d143/` (contains `output.bnd`, `output.cfg`).
  - PhysiCell config exported as:  
    `/home/mruscone/Desktop/github/mcp-biomodelling-servers/PhysiCell/artifacts/c8afd751-a294-4ae5-b5fb-e4f296bf90e9/PhysiCell\_TNF\_cancer.xml`  
    Copy this into your PhysiCell project and run (from that project): `./myproject PhysiCell\_TNF\_cancer.xml`.

At this point, the multiscale loop is closed: TNF diffuses, cells sense TNF via TNFR nodes, MaBoSS updates their fate logic, and the resulting Boolean states drive apoptosis vs proliferation in PhysiCell.

If you’d like, next I can:
- Tune TNF sources (e.g., boundary Dirichlet or secreting cells) and initial cancer cell placement, or
- Define and document a set of in silico experiments (e.g., anti-apoptotic overexpression, TNF pulses) using this model.

\end{copilotbox}
\bigskip\hrule\bigskip
\section{Request 7}
\begin{userbox}
\textbf{User Request}
\quad \textit{2026-03-04 11:11:58 UTC}
\quad Model: \texttt{copilot/gpt-5.1}

Tune TNF resources

\end{userbox}
\begin{thinkbox}
\textbf{Thinking:} \textit{Thinking}

**Interpreting TNF Resources**

The user is asking me to tune TNF resources in PhysiCell. They might want adjustments to boundary conditions, cell secretion, or pulses. I should propose a concrete setup for TNF source profiles, noting we currently have no TNF due to specific settings. They didn't specify which method to use, so I’ll suggest options like a Dirichlet boundary for a perfused source or adding TNF-secreting cells. I can choose a reasonable approach for this.

\end{thinkbox}
\begin{thinkbox}
\textbf{Thinking:} \textit{Thinking}

**Tuning TNF Resources**

The user asked me to "Tune TNF resources," suggesting I should adjust substrate parameters and interactions for a more realistic setup. I could set the initial TNF to a low baseline and add a Dirichlet boundary to mimic inflow, or implement a transient bolus, though that might need more specific fields. A better approach could be to increase the initial TNF condition, enable Dirichlet, and potentially add TNF secretion from cancer cells. The user's mention of "resources" might mean both substrate and cell interactions, so I’ll plan my adjustments accordingly.

\end{thinkbox}
\begin{thinkbox}
\textbf{Thinking:} \textit{Thinking}

**Adjusting TNF Tools**

I have tools for adding substrates and setting interactions, but there’s no clear update option. I might call `add\_single\_substrate` again with the same name to override parameters, and if that doesn't work, I can manually reconfigure by re-adding the substrate. While doing this, I’ll also adjust the cell TNF interactions to a typical uptake of 0.01, focusing on supply. Adding a second cell type for TNF secretion could also be beneficial, but I want to keep the changes minimal, like establishing a TNF Dirichlet boundary and a small basal secretion from cancer cells.

\end{thinkbox}
\begin{thinkbox}
\textbf{Thinking:} \textit{Thinking}

**Configuring TNF Substrate and Cells**

I’m planning to set up the TNF substrate with an initial condition of 0, enabling Dirichlet with a value possibly around 10 ng/mL. I’ll also add a new cell type called 'TNF\_source\_cell' that secretes TNF strongly, so it can create a significant gradient nearby. However, since the user didn’t mention needing additional cell types, I’ll focus on treating TNF as diffusing only from the environment. I’ll check the simulation summary and re-add the TNF substrate with the updated settings.

\end{thinkbox}
\begin{thinkbox}
\textbf{Thinking:} \textit{Reviewed TNF resources and adjusted related settings}

**Adjusting Cancer Cell Uptake and TNF Settings**

I’m considering adjusting the cancer cell uptake rate to possibly 0.05 for longer TNF persistence, but since "tuning resources" likely focuses on supply, I might lower the uptake from 0.1 to 0.01 instead. I can call `set\_substrate\_interaction` again with new rates to overwrite the existing settings.

Next, I’ll use `mcp\_physicell\_add\_single\_substrate` for the updated TNF settings and adjust the cancer cell uptake. After that, I’ll export the modified XML under a new filename and explain how the TNF profile will look while keeping the instructions concise. Let's get started with the tools!

\end{thinkbox}
\begin{copilotbox}
\textbf{GitHub Copilot:}

I’ll adjust the TNF substrate source and cell uptake, then regenerate the PhysiCell XML with tuned TNF supply.

\end{copilotbox}
\begin{toolokbox}
\textbf{Tool:} \texttt{mcp\_physicell\_add\_single\_substrate} (\checkmark)
\textit{Running `add\_single\_substrate`}
\begin{lstlisting}
**Substrate added:** TNF
- Diffusion: 50000 μm²/min
- Decay: 0.01 min⁻¹
- Initial: 0 ng/mL
- Dirichlet boundaries: xmin, xmax, ymin, ymax = 10 ng/mL
- Progress: 80%
**Next step:** Use `add_single_cell_type()` to add cancer cells, immune cells, etc.
{
  "result": "**Substrate added:** TNF\n- Diffusion: 50000 μm²/min\n- Decay: 0.01 min⁻¹\n- Initial: 0 ng/mL\n- Dirichlet boundaries: xmin, xmax, ymin, ymax = 10 ng/mL\n- Progress: 80%\n**Next step:** Use `add_single_cell_type()` to add cancer cells, immune cells, etc."
}
\end{lstlisting}
\end{toolokbox}
\begin{toolokbox}
\textbf{Tool:} \texttt{mcp\_physicell\_set\_substrate\_interaction} (\checkmark)
\textit{Running `set\_substrate\_interaction`}
\begin{lstlisting}
**Substrate interaction set:** cancer_cell ↔ TNF (secretion: 0, uptake: 0.02 min⁻¹)
{
  "result": "**Substrate interaction set:** cancer_cell ↔ TNF (secretion: 0, uptake: 0.02 min⁻¹)"
}
\end{lstlisting}
\end{toolokbox}
\begin{copilotbox}
\textbf{GitHub Copilot:}

I’ve added boundary TNF sources and retuned cell uptake; next I’ll export an updated XML so you can run with these TNF conditions.

\end{copilotbox}
\begin{toolokbox}
\textbf{Tool:} \texttt{mcp\_physicell\_export\_xml\_configuration} (\checkmark)
\textit{Running `export\_xml\_configuration`}
\begin{lstlisting}
## XML Configuration Exported

**File:** /home/mruscone/Desktop/github/mcp-biomodelling-servers/PhysiCell/artifacts/c8afd751-a294-4ae5-b5fb-e4f296bf90e9/PhysiCell_TNF_cancer_tunedTNF.xml (14KB)
**Source:** Created from scratch
**Substrates:** 1 (TNF)
**Cell Types:** 1 (cancer_cell)
**Progress:** 80%

**Next step:** Copy to PhysiCell project directory and run:
```bash
./myproject PhysiCell_TNF_cancer_tunedTNF.xml
```
{
  "result": "## XML Configuration Exported\n\n**File:** /home/mruscone/Desktop/github/mcp-biomodelling-servers/PhysiCell/artifacts/c8afd751-a294-4ae5-b5fb-e4f296bf90e9/PhysiCell_TNF_cancer_tunedTNF.xml (14KB)\n**Source:** Created from scratch\n**Substrates:** 1 (TNF)\n**Cell Types:** 1 (cancer_cell)\n**Progress:** 80%\n\n**Next step:** Copy to PhysiCell project directory and run:\n```bash\n./myproject PhysiCell_TNF_cancer_tunedTNF.xml\n```"
}
\end{lstlisting}
\end{toolokbox}
\begin{copilotbox}
\textbf{GitHub Copilot:}

Here is how I’ve tuned the TNF resources and what it implies:

- **TNF supply (environment):**
  - Enabled **Dirichlet boundaries** for TNF at all domain edges with a fixed value of **10 ng/mL**.
  - TNF still has diffusion coefficient **50,000 μm²/min** and decay **0.01 min⁻¹**, so you now get a sustained inflow from the tissue boundaries and a gradient toward the center.

- **TNF consumption (cells):**
  - Updated **cancer\_cell–TNF interaction** to `secretion\_rate = 0`, `uptake\_rate = 0.02 min⁻¹` (less aggressive than 0.1), so cells consume TNF but don’t instantly deplete it, allowing a more persistent TNF field that still depends on distance from boundaries and local cell density.

- **Updated configuration file:**
  - New XML with these TNF settings:

\end{copilotbox}
\begin{copilotbox}
\textbf{GitHub Copilot:}

- Run from your PhysiCell project as:  
    `./myproject PhysiCell\_TNF\_cancer\_tunedTNF.xml`

This setup gives you a controlled, constant TNF “bath” from the boundaries while letting cell density and uptake sculpt local TNF, which then feeds into the MaBoSS TNFR nodes and fate decisions. If you’d like, I can next define specific initial cancer cell placements or design a TNF pulse protocol instead of constant boundaries.

\end{copilotbox}
\begin{copilotbox}
\textbf{GitHub Copilot:}

I’ll adjust the TNF substrate source and cell uptake, then regenerate the PhysiCell XML with tuned TNF supply.

\end{copilotbox}
\begin{toolokbox}
\textbf{Tool:} \texttt{mcp\_physicell\_add\_single\_substrate} (\checkmark)
\textit{Calling `mcp\_physicell\_add\_single\_substrate`}
\end{toolokbox}
\begin{toolokbox}
\textbf{Tool:} \texttt{mcp\_physicell\_set\_substrate\_interaction} (\checkmark)
\textit{Calling `mcp\_physicell\_set\_substrate\_interaction`}
\end{toolokbox}
\begin{copilotbox}
\textbf{GitHub Copilot:}

I’ve added boundary TNF sources and retuned cell uptake; next I’ll export an updated XML so you can run with these TNF conditions.

\end{copilotbox}
\begin{toolokbox}
\textbf{Tool:} \texttt{mcp\_physicell\_export\_xml\_configuration} (\checkmark)
\textit{Calling `mcp\_physicell\_export\_xml\_configuration`}
\end{toolokbox}
\begin{copilotbox}
\textbf{GitHub Copilot:}

Here is how I’ve tuned the TNF resources and what it implies:

- **TNF supply (environment):**
  - Enabled **Dirichlet boundaries** for TNF at all domain edges with a fixed value of **10 ng/mL**.
  - TNF still has diffusion coefficient **50,000 μm²/min** and decay **0.01 min⁻¹**, so you now get a sustained inflow from the tissue boundaries and a gradient toward the center.

- **TNF consumption (cells):**
  - Updated **cancer\_cell–TNF interaction** to `secretion\_rate = 0`, `uptake\_rate = 0.02 min⁻¹` (less aggressive than 0.1), so cells consume TNF but don’t instantly deplete it, allowing a more persistent TNF field that still depends on distance from boundaries and local cell density.

- **Updated configuration file:**
  - New XML with these TNF settings:  
    /home/mruscone/Desktop/github/mcp-biomodelling-servers/PhysiCell/artifacts/c8afd751-a294-4ae5-b5fb-e4f296bf90e9/PhysiCell\_TNF\_cancer\_tunedTNF.xml  
  - Run from your PhysiCell project as:  
    `./myproject PhysiCell\_TNF\_cancer\_tunedTNF.xml`

This setup gives you a controlled, constant TNF “bath” from the boundaries while letting cell density and uptake sculpt local TNF, which then feeds into the MaBoSS TNFR nodes and fate decisions. If you’d like, I can next define specific initial cancer cell placements or design a TNF pulse protocol instead of constant boundaries.

\end{copilotbox}
\end{document}